\section{Finite Forcing Conditions}

Once counting has wholes and parts, the grammar can approach a real
number without possessing a completed continuum. It does so by finite
conditions. A finite condition is a record that has gone far enough to
constrain the next readings, but not so far that it claims to have
settled every possible refinement. It is a disciplined partial answer.
It permits extension; it forbids the pretense that extension has already
ended.

This is the point at which residue enters. A residue is not a defect in
the record. It is the record's honest remainder: the direction and size
of what has not yet been absorbed. When an approximation improves, the
residue shrinks or changes shape. The shrinking is readable because the
earlier count has taught the grammar how to preserve wholes and expose
parts. The residue carries the part that remains after the current whole
has done all it can do.

The real number, in this chapter's foundation, is therefore not an
infinite treasure waiting at the end of a completed road. It is the
stable reading forced by a finite approach whose residues are bounded
and driven downward. The grammar may say: every admissible extension is
being pressed toward this reading. It may not say: all extensions have
already been traversed. The difference is not a technical caution. It is
the grammatical difference between measurement and omniscience.

The familiar decimal image helps if it is used carefully. A string such
as 0.999... can be read as a completion, but the grammar at the
foundation reads instead the finite discipline of approach: 0.9, then
0.99, then 0.999, with each stage carrying a residue that has become
smaller. The equality with the limiting value is not an initial gift. It
is the reading made available when the residue has been bounded in the
right way. The grammar earns the identification through extension; it
does not assume it before extension begins.

This stance also clarifies the continuum problem. A continuous model is
tempting because it promises every intermediate distinction at once. But
the grammar of measurement cannot accept unrecorded distinctions merely
because a completed continuum would be elegant. Any intermediate
structure that matters must be recoverable from finite refinement. If it
cannot be recovered from a ledger of admissible readings, it may be a
useful idealization, but it is not a primitive measurement.

The Cantor--G\"odel--Cohen illustration belongs here. The lesson for this
chapter is not that set theory has failed to decide a famous sentence.
The lesson is that completions can differ while agreeing on all the
finite records the grammar has chosen to carry. A completion is a way of
reading the refinement history. It is not automatically a new
measurement. The grammar permits completion only when the added
structure remains answerable to recoverable refinement.

Finite forcing conditions thus give the book its first account of
approach. A finite condition decides what it can decide. It leaves a
residue where it must leave one. It narrows future readings without
pretending to replace them. It carries the past forward as constraint,
not as a total possession of the future.

The next obligation is temporal. A finite approach needs stages, and
stages need a grammar of before and after. But time cannot be imported as
an already understood background. The grammar has to obtain time from
the same source as number: a permitted alternation that can be read
again. The clock enters as the finite condition's way of becoming a
trial.
